\documentclass[11pt]{article}
\usepackage[T1]{fontenc}
\usepackage{lmodern,amsmath,amssymb,amsthm,booktabs}
\usepackage[hmargin=1.15in,vmargin=1.10in]{geometry}
\usepackage[hidelinks]{hyperref}
\input{glyphtounicode}
\pdfgentounicode=1
\newtheorem{theorem}{Theorem}[section]
\newtheorem{lemma}[theorem]{Lemma}
\newtheorem{proposition}[theorem]{Proposition}
\newtheorem{corollary}[theorem]{Corollary}
\newtheorem{remark}[theorem]{Remark}
\DeclareMathOperator{\Li}{Li}
\title{A zero estimate for triangular logarithmic systems\\
and residue discs in Chabauty--Kim loci}
\author{Prajay Basu}
\date{}
\begin{document}
\maketitle
\begin{abstract}
For a strictly triangular logarithmic system of rank \(N\) in characteristic
\(p\) with a one-dimensional common left kernel, normalise the solution
on that kernel. We bound the vanishing order of a polynomial linear form of degree \(A\)
by \(NA+d\binom N2\) whenever this quantity is below \(p\).
Here \(d\) depends on the poles and the centre of expansion.
A family with two fixed poles and constant linear forms has quadratic
vanishing order in the rank. Applied to finite polylogarithms, the estimate
bounds the number of residue discs occupied by the refined depth-\(2s\)
Chabauty--Kim locus of the thrice-punctured line by
\(O_s(p^{1-1/(4s)})\), where \(s\) is the number of inverted primes.
For \(S=\{2,3\}\) we identify a Bernoulli factor in a depth-four cofactor
and determine the refined locus at the auxiliary primes
\(16843\), \(1006003\), and \(2124679\).
\end{abstract}

\noindent\textbf{2020 Mathematics Subject Classification.}
Primary 12H20; Secondary 11G30, 11G55.

\section{Introduction}
A single equation between Coleman functions can have zeros in many residue
discs. To bound the number of discs meeting a Chabauty--Kim locus, we
eliminate the coordinate \emph{within} a disc. What remains is an equation
over \(\mathbb F_p\) between the residue and its finite polylogarithms.
Its coefficients can vary with the prime. The estimate must therefore
apply to every nonzero equation of the prescribed degrees.

Cyclic differentiation and minors also enter characteristic-zero
multiplicity estimates \cite{Fischler}. Here the height bound on a
normalised row controls the extra derivative constants in
\(\mathbb F_p(X^p)\).

\begin{theorem}[Triangular logarithmic zero estimate]\label{thm:zero}
Let \(k=\mathbb F_p\), let \(\Sigma\subset k\) be finite and nonempty, and
let \(b\in k\). In a common basis let \(A_a\) for \(a\in\Sigma\) be strictly
upper triangular constant \(N\times N\) matrices whose common left kernel
is spanned by a coordinate row \(e\). Write
\(T=\prod_{a\in\Sigma\cup\{b\}}(X-a)\) and \(d=\deg T-1\).
Suppose a polynomial column \(w\) satisfies \(ew=1\) and
\[
 T w'\equiv\sum_{a\in\Sigma}\frac{T}{X-a} A_a w\pmod{(X-b)^p}.
\]
If \(P\ne0\) is a polynomial row of degree at most \(A\) and
\(H=NA+d\binom{N}{2}<p\), then \(\operatorname{ord}_b(Pw)\le H\).
For two fixed poles and an ordinary centre there are systems of unbounded
rank \(R\) and constant rows whose evaluations have order at least
\(R^2/4-1\) and are nonzero modulo \((X-b)^p\), for every sufficiently
large prime for each system.
\end{theorem}

Let \(X=\mathbb P^1\setminus\{0,1,\infty\}\) with its standard integral
model. A residue disc is \emph{occupied} by a subset of \(X(\mathbb Z_p)\)
if it contains a point of that subset. Thus the quantity we bound is the
cardinality of its image in \(X(\mathbb F_p)\). We use the refined and
unrefined Selmer conditions and the lower central series depth
of \cite{BKL}.
\begin{theorem}[Occupied residue discs]\label{thm:discs}
For any finite set \(S\) of \(s\ge1\) rational primes and prime \(p\notin S\),
the full refined Chabauty--Kim locus of \(X\) at depth \(2s\) occupies
\(O_s(p^{1-1/(4s)})\) discs. The full unrefined locus at depth \(4s\)
occupies \(O_s(p^{1-1/(8s)})\) discs. The constants depend only on \(s\).
In particular the refined exponents at depths four for \(s=2\) and two
for \(s=1\) are \(7/8\) and \(3/4\). For the inverse image of \(X\) in
\(y^2=x^5-x\) under the hyperelliptic projection, the refined depth-two
bound for \(S=\{2\}\) is also \(O(p^{3/4})\).
\end{theorem}

At depth four, L\"udtke's equations for \(S=\{2,3\}\) use two known
points \(3,9\) to eliminate the global period coefficients of a
determinant \cite[Theorem~3.8]{Ludtke}. We use \(\log_p(p)=0\),
\(Y_m(x)=p^{-m}\Li_m(x)\), and \(k=\log_p(2)/p\). Define \(A\) as the
first signed cofactor of the two anchor rows
\[
 (Y_4(x),\,jY_3(x),\,j^3Y_1(x)),\qquad (x,j)=(3,1),(9,2).
\]
This fixes the sign and the powers of the logarithm removed from the
determinant.
\begin{theorem}[Leading cofactor]\label{thm:cofactor}
For every prime \(p\ge7\),
\(9A\equiv52kB_{p-3}\pmod p\), where \(B_j\) is the \(j\)-th Bernoulli
number. For \(p\ne13\), the leading digit vanishes exactly when
\(2^{p-1}\equiv1\pmod{p^2}\) or \(p\) divides the numerator of \(B_{p-3}\).
\end{theorem}

\begin{theorem}[Depth four and small torsion orders]\label{thm:cert}
For \(S=\{2,3\}\) the full refined depth-four locus equals
\(X(\mathbb Z[1/6])\) at each of the six primes
\[
 11,\quad1093,\quad3511,\quad16843,\quad1006003,\quad2124679.
\]
On the \((1,1)\) refinement, weight two alone excludes primitive fourth
roots of unity for \(7\le p<10^7\), and primitive third and sixth roots for
\(7\le p<2\cdot10^7\), whenever the roots lie in \(\mathbb Q_p\).
\end{theorem}

The first three primes of Theorem~\ref{thm:cert} are in L\"udtke's
theorem for \(5\le p<10000\) \cite[Theorem~1.2]{Ludtke}.
The remaining three calculations use all-disc finite-field screens and
local precision bounds. The torsion ranges use the published searches of
Me\v{s}trovi\'c and Sun \cite{Mestrovic,Sun}.
Li and L\"udtke previously used finite polylogarithms to constrain
Teichm\"uller coordinates over imaginary quadratic fields at infinite
depth \cite[Theorem~8.26]{LL}. Our count uses a fixed list of weights over
\(\mathbb Q\). Zehavi studies eventual dimension inequalities for
projective curves at sufficiently large depth \cite{Zehavi}.
For the punctured line, the dimension count produces a relation at the
depths in Theorem~\ref{thm:discs}; bounding the roots of its reduction
is the additional step needed for a power saving in \(p\).

\section{Multiplicity at a logarithmic jet}\label{sec:zero}

We work over $k=\mathbb F_p$. Let $\Sigma\subset k$ be a nonempty set of
distinct finite poles, $b\in k$ a centre, and $(A_a)_{a\in \Sigma}$ constant
strictly upper triangular $N\times N$ matrices in one ordered basis.
Put $T=\prod_{a\in \Sigma\cup\{b\}}(X-a)$ and $d=\deg T-1$. Suppose the
common left kernel of the $A_a$ is generated by a coordinate row $e$.
The polynomial column $w$ in Theorem~\ref{thm:zero} has $ew=1$ and satisfies
\begin{equation}\label{eq:log-jet}
 T w'\equiv\sum_{a\in \Sigma}\frac{T}{X-a}A_aw\pmod{(X-b)^p}.
\end{equation}
At an ordinary centre $b\notin \Sigma$ this jet exists for every initial
column $w_0$ with $ew_0=1$. Indeed, expand $\sum_{a\in\Sigma} A_a/(X-a)$ at $b$
and solve for the coefficients of $w$ through degree $p-1$. Computing
coefficient $j$ requires division by $j$ for $1\le j<p$. The initial
condition and $eA_a=0$ give $ew=1$ for the truncated polynomial.
Multiplication by $T$ raises the order of the remaining differential
defect from at least $p-1$ to at least $p$. At a singular centre
we use \eqref{eq:log-jet} as a hypothesis.

\begin{proof}[Proof of Theorem~\ref{thm:zero}]
On polynomial rows define
\[
 D h=T h'+\sum_{a\in \Sigma}\frac{T}{X-a}h A_a .
\]
If $P$ is a nonzero row of degree at most $A$, then $D^jP$ has entries
of degree at most $A+jd$ by induction. Equation~\eqref{eq:log-jet}
gives $(Dh)w\equiv T(hw)'\pmod{(X-b)^p}$. Since $X-b$ divides $T$,
the operator $T\,d/dX$ preserves $(X-b)^m$ for every $m\ge0$.
Consequently $Pw\equiv0\pmod{(X-b)^m}$ implies
$(D^jP)w\equiv0\pmod{(X-b)^m}$ for all $j$ whenever $m\le p$.

Over $K=k(X)$ take the first $r$ successive independent rows
$P,DP,\ldots,D^{r-1}P$, stopping when the next row is dependent.
Such an $r$ exists with $1\le r\le N$. Write $M$ for their matrix.
Its row span $V$ is stable under $D$, since $D(fh)=Tf'h+fDh$,
and hence also under $\nabla=T^{-1}D$.
Let $I$ be the ordered pivot columns of $M$ and let $B=M_I$.
The rows of $B^{-1}M$ form the reduced echelon basis of $V$.
Its last row $h$ starts with a $1$ in the last pivot column.
In $\nabla h=h'+\sum_{a\in\Sigma} hA_a/(X-a)$, all coordinates through this
column vanish by strict upper triangularity. Every nonzero element
of $V$ starts in a pivot column, so $\nabla h=0$.

All $r$-column minors of $M$ have degree at most
$H_r=rA+d\binom r2$: each determinant uses one entry from each
successive derivative row. By Cramer's rule each coordinate of $h$
is a ratio of such a minor (possibly zero) to $q=\det B\ne0$.
In particular both degrees are bounded by $H_r\le H=NA+d\binom N2<p$.

We now show in increasing order of $j$ that $h_j\in k$.
Horizontality gives
\[
 h_j'=-\sum_{a\in \Sigma}\frac{1}{X-a}\sum_{i<j}h_i(A_a)_{ij}.
\]
By induction all numerators on the right are in $k$.
The formal Laurent derivative at $a$ has zero coefficient at
$(X-a)^{-1}$, since only the derivative of the constant term could
contribute. Taking residues at the distinct $a$ makes every numerator
zero. Hence $h_j'=0$. To see why this implies $h_j\in k$ here, write a
nonzero coordinate in lowest terms as $u/v$. The identity $u'v=uv'$
gives $v\mid v'$, whence $v'=0$ by degree and then $u'=0$.
Both polynomials have degree less than $p$, so coefficient comparison
makes both constant. The zero coordinate is already constant.
This proves the induction, including the empty sum for the first
coordinate. The same residue equations give $hA_a=0$ for every $a$.
The common-kernel hypothesis and the pivot entry $1$ give $h=e$.

The identity $q B^{-1}=\operatorname{adj}(B)$ expresses $qe$ as a
polynomial linear combination of the rows of $M$. If $Pw$ vanished
to order at least $H_r+1$, the propagation of vanishing would apply
at this precision because $H_r+1\le p$. It would give
$(X-b)^{H_r+1}\mid qew=q$, contradicting $0\ne q$ and $\deg q\le H_r$.
Thus $Pw$ has a nonzero coefficient of order at most $H_r\le H$.
\end{proof}

\subsection{Classical finite polylogarithms}
For $1\le i\le n$ put $\ell_i(X)=\sum_{m=1}^{p-1}X^m/m^i$.
The following form of Theorem~\ref{thm:zero} will be used in
Section~\ref{sec:discs}.

\begin{corollary}\label{cor:box}
Let $n\ge1$, $t,A\ge0$, $N=\binom{n+t}{n}$ and
$H=NA+\binom N2$. If $t<p$ and $H<p$, every nonzero
$f\in k[X,y_1,\ldots,y_n]$ of $X$-degree at most $A$ and total
$y$-degree at most $t$ has
\[
 \operatorname{ord}_0 f(X,\ell_1(X),\ldots,\ell_n(X))\le H.
\]
In particular substitution modulo $X^p$ is injective on this degree box.
\end{corollary}

\begin{proof}
Use the $N$ monomials of total degree at most $t$ as coordinates,
ordered by decreasing weight for $\mathrm{wt}(y_i)=i$.
The two derivations
\[
 E=\partial_{y_1},\qquad L=\sum_{i=2}^n y_{i-1}\partial_{y_i}
\]
act on this space by constant strictly upper triangular row matrices:
each lowers weight by one. If $Ef=Lf=0$ on this monomial space, the
identities $[\partial_{y_j},L]=\partial_{y_{j+1}}$ for $1\le j<n$
successively give $\partial_{y_j}f=0$ for every $j$. Each exponent
is below $p$, so $f$ is constant. The common left kernel is precisely
the coordinate line of the monomial $1$.

Take $w$ to be the column obtained by substituting $y_i=\ell_i$ in the
monomials. Differentiating the sums gives
$X(1-X)\ell_1'=X-X^p$ and $X\ell_i'=\ell_{i-1}$ for $i\ge2$.
The product rule now gives \eqref{eq:log-jet} with $b=0$ and residues
$L$ at $0$ and $-E$ at $1$: the connection on monomials is
$L/X+E/(1-X)$. Here $d=1$. Apply the theorem to the row of
coefficients of $f$ to obtain the claimed inequality.
\end{proof}

\subsection{Words and strict multiple polylogarithms}
Take a finite set of distinct words on an alphabet $\Sigma\subset k$,
containing the empty word and closed under deletion of the first letter.
Order the words by decreasing length. The residue for $a$ sends the row
of $av$ to the row of $v$ and kills rows with any other first letter,
including the empty row. These matrices are strictly upper triangular.
In the image under the $a$-matrix, the $v$-coordinate recovers the
original $av$-coefficient whenever $av$ occurs. Thus the common left
kernel is the empty-word line. Multiplying individual residue matrices
by nonzero constants preserves this calculation. Theorem~\ref{thm:zero}
applies to polynomial-coefficient linear forms on these word jets.

For example, with $s_j\ge1$, consider
\[
 F_{s_1,\ldots,s_h}(X)=
 \sum_{p>n_1>\cdots>n_h\ge1}
 \frac{X^{n_1}}{n_1^{s_1}\cdots n_h^{s_h}},\qquad F_{\varnothing}=1 .
\]
Encode this index by $0^{s_1-1}1\cdots0^{s_h-1}1$ and take all suffixes
of the words in question. If $s_1>1$, then
$XF'_{s_1,\mathbf t}=F_{s_1-1,\mathbf t}$. For $s_1=1$ summing the
geometric progression $\sum_{n=m+1}^{p-1}X^n$ above the largest tail
index $m$ (take $m=0$ for the empty tail) gives
\[
 X(1-X)F'_{1,\mathbf t}=XF_{\mathbf t}(X)-X^p F_{\mathbf t}(1).
\]
Modulo $X^p$ the residues are therefore the deletion operators for $0$
and $1$, with a minus sign at $1$. The estimate holds at $b=0$ for
linear forms in distinct words. For a polynomial expression a separate
check of its expansion is needed. For instance $F_1^2-2F_{1,1}$ vanishes
modulo $X^p$ for odd $p$: at each $1\le n<p$ its two terms agree by
$1/(i(n-i))=(1/i+1/(n-i))/n$ summed over $1\le i<n$.

\subsection{Quadratic growth with two fixed poles}
For every $n\ge2$ we construct a system of rank $R$ with
$n+1\le R\le 2n$ and a constant row whose order at an ordinary centre
is at least $n^2-1$. Work first over $\mathbb Q[[X]]$ and define
\[
 J_u f=\int_0^X\frac{f(t)}{1-t}\,dt,\qquad
 J_v f=\int_0^X\frac{f(t)}{1+t}\,dt .
\]
The first $n^2-1$ Taylor coefficients (including the constant) of
\[
 G=\sum_{0\le i,j<n} b_{ij}J_u^iJ_vJ_u^j1
\]
give a rational linear map from dimension $n^2$ to dimension $n^2-1$.
Choose a nonzero kernel element, clear its denominators, and let $k_0$
be the last index of a nonzero row. Introduce the ordered coordinates
$U_0,\ldots,U_{k_0},V_0,\ldots,V_{n-1}$ by the recursions
\[
 V_{n-1}=1,\qquad V_j=J_uV_{j+1}\ (j<n-1),\qquad
 U_i=J_uU_{i+1}+J_v\sum_{j=0}^{n-1}b_{ij}V_{n-1-j},\qquad U_{k_0+1}=0,
\]
where $U_{k_0+1}$ is a boundary value. Expanding recursively gives
$U_0=G$. Differentiation gives a strictly upper triangular system with
poles at $1,-1$ and rank $R=k_0+1+n$. Up to sign the residue at $1$
shifts each of the two chains to the next coordinate. Its left kernel
has basis the terminal rows $e_{U_{k_0}},e_{V_{n-1}}$. The residue at
$-1$ kills $e_{V_{n-1}}$ and takes $e_{U_{k_0}}$ to the nonzero
reversed last array row in the $V$-coordinates. The common kernel is
therefore exactly $k e_{V_{n-1}}$ after reduction at every prime for
which the last array row remains nonzero.

At the ordinary centre $0$ the Taylor coefficients through degree $p-1$
are $p$-integral: the kernels $1/(1\pm X)$ have integral coefficients
and each integration divides the coefficient of degree $j-1$ by $j$.
Reduce this finite jet for an odd prime $p>\max\{n^2,2\binom R2\}$ at
which the last array row survives. Its $U_0$-coordinate has vanishing
order at least $n^2-1$ by construction. Theorem~\ref{thm:zero}, with
$d=2$ and $A=0$, shows that this coordinate is nonzero modulo $X^p$.
As $n$ varies, $R$ is unbounded and $n^2-1\ge R^2/4-1$. Consequently
even two fixed poles and a constant input row admit quadratic growth
in rank. Any improvement to linear growth for the monomial systems of
Corollary~\ref{cor:box} would require a property specific to that family.

\section{Residue discs at fixed depth}\label{sec:discs}

The proof of Theorem~\ref{thm:discs} uses a relation whose coefficients
may depend on every $p$-adic period occurring in localisation. Its degree
will depend only on $s=|S|$. Ghosh and Ward obtained
$O_n(p/\log p)$ for the roots of a single $\ell_n$ with $n\ge2$
\cite[Theorem~1]{GW}. In the Chabauty--Kim setting,
Betts, Kumpitsch and L\"udtke used a finite polylogarithm of weight
$p-3$ to isolate one Teichm\"uller point
\cite[Proposition~5.15]{BKL}. Here the weights stay fixed and the
relation may mix them. We first bound the number of residues satisfying
any such relation. Write
\[
  \ell_i(X)=\sum_{m=1}^{p-1}\frac{X^m}{m^i}\in\mathbb F_p[X],
  \qquad \boldsymbol\ell=(\ell_1,\ldots,\ell_n).
\]

\begin{lemma}\label{lem:mixed-disc-count}
Fix $n\ge2$ and nonnegative integers $h,d$. For every prime $p$ and every
nonzero $R\in\mathbb F_p[X,y_1,\ldots,y_n]$ with $\deg_X R\le h$
and total $y$-degree at most $d$,
\[
 \#\{a\in\mathbb F_p\setminus\{0,1\}:
                  R(a,\boldsymbol\ell(a))=0\}
       \ll_{n,h,d}p^{1-1/(2n)} .
\]
\end{lemma}

\begin{proof}
We use the auxiliary variables $X,W,y_1,\ldots,y_n$ in the manner of
Ghosh and Ward \cite[\S3]{GW}. Put
\[
 \delta=X(1-X)\partial_X+(X-W)\partial_{y_1}
             +(1-X)\sum_{i=2}^n y_{i-1}\partial_{y_i},
 \qquad \delta W=0.
\]
Upon substituting $W=X^p$ and $y_i=\ell_i(X)$ this derivation induces
$D=X(1-X)\,d/dX$. Indeed $X(1-X)\ell_1'=X-X^p$ by summing a geometric
progression, and $X\ell_i'=\ell_{i-1}$ for $i\ge2$.

Set $N_t=\binom{t+n}{n}$, choose an integer $b>3nd+h+1$, and put
$A=r=t^2$, $B=bt$. Seek $\Phi$ in the space of polynomials of
degrees at most $A,B,t$ in $X,W,\boldsymbol y$, respectively, taking
total degree in $\boldsymbol y$. For $0\le j<r$ impose the linear conditions
\[
             R\mid (\delta^j\Phi)(X,X,\boldsymbol y).
\]
The right-hand polynomials have $X$-degree at most $E=A+B+r$ and
$y$-degree at most $t$, since $\delta$ raises total $(X,W)$-degree by
at most one and preserves the $y$-degree bound. In this space, multiples
of $R$ include the injective image under multiplication by $R$ of the
box with bounds $E-h,t-d$. For $t$ sufficiently large the codimension is
therefore at most
\[
 (E+1)(N_t-N_{t-d})+hN_{t-d}
       \le (3ndt+h)N_t .
\]
Here $E+1\le 3t^2$ and telescoping the last $d$ binomial increments gives
$N_t-N_{t-d}\le d\binom{t+n-1}{n-1}\le dnN_t/(t+1)$.
The total codimension is smaller than $(A+1)(B+1)N_t$, so a nonzero
$\Phi$ exists.

Choose $t=\lfloor c_n p^{1/(2n)}\rfloor$ with $c_n>0$ sufficiently small.
For $A=t^2$, the bound $N_t A+\binom{N_t}{2}=O_n(t^{2n})$ uses $n\ge2$.
Thus for large $p$ the bound is strictly below $p$, and $t,r<p$.
Corollary~\ref{cor:box} now shows that for every nonzero
coefficient $f(X,\boldsymbol y)$ of $\Phi$ as a polynomial in $W$,
\[
                       X^p\nmid f(X,\boldsymbol\ell(X)).
\]
In particular, taking the first nonzero coefficient, with index $i$,
the corresponding summand of
$F(X)=\Phi(X,X^p,\boldsymbol\ell(X))$ has order strictly below $(i+1)p$
at zero, and all subsequent summands have order at least $(i+1)p$.
Hence $F\ne0$ and $\deg F\le A+pB+(p-1)t$.

At a residue $a$ counted by the lemma, $a^p=a$, so the divisibility
conditions give $(D^jF)(a)=0$ for all $j<r$. If $F$ had order $m<r$
there, with leading coefficient $u$ in $X-a$, we would have
$(D^mF)(a)=m![a(1-a)]^mu\ne0$. Thus each residue contributes at
least $r$ to the degree of $F$, giving an upper bound
$1+(b+1)p/t$. The finitely many smaller primes can be absorbed into the
implied constant.
\end{proof}

\begin{proof}[Proof of Theorem~\ref{thm:discs}]
Use the standard integral model
$X=\mathbb P^1_{\mathbb Z}\setminus\{0,1,\infty\}$ and its usual
tangential base point at zero directed towards one for polylogarithmic
coordinates. Take $p\notin S$ and $s\ge1$. We may suppose $p>n+1$,
where $n=2s$ for the refined locus and $n=4s$ for the unrefined locus.
For smaller $p$ the trivial bound $p-2$ suffices, with a constant
depending only on $s$. Passing from the full depth-$n$ fundamental
group to its polylogarithmic quotient carries the global and local Kummer
conditions into the corresponding quotient conditions. It suffices to bound
the reductions of the resulting polylogarithmic locus.

The localisation formula of Betts, Kumpitsch and L\"udtke
\cite[Propositions~5.3 and~5.9, Theorem~5.6]{BKL}
has global coordinates $x_\ell,y_\ell$ for $\ell\in S$, of weight one,
and $z_j$ for odd $3\le j<n$, of weight $j$ (here $n$ is even).
In the refined case $2\notin S$ forces emptiness by the local
condition $X(\mathbb Z_2)=\varnothing$. If $2\in S$, there are $3^s$
refinement conditions, choosing for each $\ell$ one of the hyperplanes
$x_\ell=0$, $y_\ell=0$, $x_\ell+y_\ell=0$
\cite[Proposition~5.11]{BKL}. On each resulting linear subspace there
are $s+(s-1)=n-1$ free coordinates. In the unrefined case the whole
source has $2s+(2s-1)=n-1$ coordinates. Denote either list by $v$.
In these coordinates the local functions
\[
 L=p^{-1}\log_p,\qquad Y_m=p^{-m}\operatorname{Li}_m\quad(1\le m\le n)
\]
pull back to polynomials $L(v),Y_m(v)$ over $K=\mathbb Q_p$ of ordinary
degrees at most $1,m$, respectively. In fact each term in the
localisation formula of weight $m$ is a product of $m$ weight-one
coordinates, or of a coordinate $z_j$ and $m-j$ weight-one coordinates;
restricting to the chosen linear subspace preserves the degree bound.

Adjoin a free variable $x$. Define $W_0=-x$, $W_m=(x-1)Y_m(v)$ for $m>0$,
and form the map from $\mathbb A_K^n$ to $\mathbb A_K^{n+1}$ with coordinates
\[
 x,\qquad u_m=\sum_{j=0}^{m}W_{m-j}\frac{(-L(v))^j}{j!}
           \quad (1\le m\le n).
\]
All coordinates have degree at most $D=n+1$. Set $H=((n+1)D+1)^n$.
Polynomials of individual variable degrees at most $H$ in the target
have dimension $(H+1)^{n+1}$. After substitution their individual
variable degrees in the source are at most $(n+1)DH$. The source
space has dimension at most
\[
 ((n+1)DH+1)^n\le H(H+1)^n<(H+1)^{n+1}.
\]
Consequently a nonzero polynomial $P$ of individual degrees at most $H$
vanishes identically after substitution. Divide by a coefficient of
maximal $p$-adic absolute value to make $P$ integral with nonzero
reduction $R$. This normalization uses only coefficients of $P$ and
requires no integrality of $v$ or of the periods. Moreover,
$P(x,\boldsymbol u)$, viewed as a polynomial in $x$ and the local
coordinates $L,\boldsymbol Y$, vanishes on the product of the $x$-line
with the Zariski closure of the localisation image. This follows by
taking its coefficients in $x$, each of which pulls back to zero.

We compute $\boldsymbol u$ modulo $p$ at a local point $z$ reducing
to $a\ne0,1$. Let $\zeta$ be the Teichm\"uller lift of $a$, and write
$z=\zeta(1+pw)$ with $w\in\mathbb Z_p$. Then
$L(z)\equiv w\pmod p$. At $\zeta$, the $p$-depleted polylogarithm
$\operatorname{Li}_m(\zeta)-p^{-m}\operatorname{Li}_m(\zeta^p)$
equals $(p^m-1)Y_m(\zeta)$. Its reduction is
$\ell_m(a)/(1-a)$ by Besser \cite[Proposition~2.1]{Besser}. Hence
$Y_m(\zeta)$ is integral with reduction $\ell_m(a)/(a-1)$.
Set $Y_0(\zeta)=\zeta/(1-\zeta)$. Besser's translation formula
\cite[Proposition~2.3]{Besser} gives integral $Y_m(z)$ and
\[
 Y_m(z)\equiv\sum_{j=0}^m Y_{m-j}(\zeta)w^j/j!\pmod p .
\]
Take any integral $x$ lifting $a$. Since $(a-1)Y_0(\zeta)\equiv-a$,
the reductions of $W_0,W_1,\ldots,W_n$ are the coefficients through
degree $n$ of
\[
 \left(-a+\sum_{m=1}^n\ell_m(a)T^m\right)
                 \sum_{j=0}^n\frac{\bar w^j T^j}{j!}
                \quad\text{in }\mathbb F_p[T]/(T^{n+1}).
\]
The sum defining $u_m$ multiplies this expression by the truncated
exponential of $-\bar w T$. The two exponentials multiply to $1$
in this ring, by the binomial identity and $n<p$. We obtain
$u_m\equiv\ell_m(a)\pmod p$.

If $z$ lies in the locus under consideration, its local coordinates lie
in the localisation image or its closure, so
$R(a,\boldsymbol\ell(a))=0$. Apply Lemma~\ref{lem:mixed-disc-count} with
$h=H$ and $d=nH$. Each image contributes $O_n(p^{1-1/(2n)})$ discs.
Summing over at most $3^s$ images in the refined case, and taking the
single image in the unrefined case, gives the asserted exponents
$1-1/(4s)$ and $1-1/(8s)$.
\end{proof}

\paragraph{An affine genus-two pullback.}
Let $C$ be the smooth projective model of $y^2=x^5-x$ and let $U$ be
the inverse image of $X$ under the hyperelliptic projection. The branch
divisor consists of $0,\infty$ and the four fourth roots of unity,
which are pairwise distinct in odd characteristic. Thus the double
cover has genus two by Riemann--Hurwitz and good reduction away from
$2$. Removing the three ramification points above $0,1,\infty$ gives
a finite degree-two morphism $U\to X$ over $\mathbb Z[1/2]$.
Use the compatible integral base points $(-1,0)$ and $-1$. Base-point
change on $X$ twists path torsors by the path torsor between the two
base points and identifies the corresponding local and global Kummer
conditions; thus the disc bound computed in tangential coordinates
applies with base point $-1$ as well.

For $S=\{2\}$, functoriality sends the full refined depth-two locus of
$U$ into that of $X$: the homomorphism of unipotent fundamental groups
preserves the lower central series, and its induced map on torsors
commutes with localisation and the local Kummer maps. The image of
each local Kummer closure lies in the corresponding closure for $X$,
by continuity in the Zariski topology. At an odd prime $p$, each
residue $a$ on the target has at most two lifts, since its fibre
satisfies $y^2=a^5-a$. Theorem~\ref{thm:discs} gives
$O(p^{3/4})$ occupied discs for the full refined depth-two locus
of $U$. This conclusion concerns the open curve with the specified
three punctures.

\section{Exceptional primes for the thrice-punctured line}\label{sec:exceptional}

Throughout this section $S=\{2,3\}$, $p\ge7$, and the Coleman
logarithm satisfies $\log_p(p)=0$. We use the order $z,3,9$ for
the rows of the determinant in \cite[(6.2)]{Ludtke}. Put
\[
 L=\log_p(z)/p,\quad Y_m(z)=p^{-m}\operatorname{Li}_m(z),\quad
 k=\log_p(2)/p,\quad s=\log_p(3)/p,\quad d=Y_2(3).
\]
Define $(A,B,C)$ to be the cross product of
$(Y_4(3),Y_3(3),Y_1(3))$ and $(Y_4(9),2Y_3(9),8Y_1(9))$.
The two necessary equations on refinement $(1,0)$ are
\begin{align}
 skY_2+dLY_1&=0,\label{eq:exc-f2}\\
 s^3AY_4+s^2BLY_3+CL^3Y_1&=0.\label{eq:exc-f4}
\end{align}
The first equation is the depth-two function divided by $p^4$.
Each raw determinant row has a factor $p^4$, and the second and
third entries at the anchors have the further factors $s,s^3$.
Expansion along the first row gives \eqref{eq:exc-f4} after division
by $p^{12}s$. The scalar $s$ is nonzero since $3$ has infinite
multiplicative order. We write $f_4$ for the raw determinant below.

Jaswal uses extra divisibility of an abelian logarithm at a prime chosen
for each class on a modular curve \cite{Jaswal}. In the equations above
the valuations of the logarithms at the given auxiliary prime determine
which cofactor survives reduction.

\subsection{The coefficient of the tetralogarithm}

\begin{proof}[Proof of Theorem~\ref{thm:cofactor}]
Let $Z=p^{-3}\zeta_p^{\mathrm{Col}}(3)$, in the normalization
$\operatorname{Li}_3(1)=\zeta_p^{\mathrm{Col}}(3)$.
Since $\operatorname{Li}_1(x)=-\log_p(1-x)$, the third anchor
entries are $-k$ and $-24k$. Thus
\[
 A=2k\bigl(Y_3(9)-12Y_3(3)\bigr).
\]
Duplication gives $Y_3(9)=4(Y_3(3)+Y_3(-3))$.
The trilogarithm identity
$Y_3(-3)-2Y_3(3)=-13Z/6$ follows from
\cite[\S5.7, (5.21)]{CDL}. Hence $A=-52kZ/3$.
Duplication at $1$ also gives $Y_3(-1)=-3Z/4$.
By \cite[Proposition~2.1]{Besser},
\[
 Y_3(-1)\equiv-\tfrac12\sum_{j=1}^{p-1}(-1)^j j^{-3}
 \pmod p.
\]
The full reciprocal-cube sum is zero in $\mathbb F_p$. Consequently
the alternating sum equals twice its even-index part, namely
$\tfrac14\sum_{i=1}^{(p-1)/2}i^{p-4}$. Faulhaber's formula reduces
the latter sum, before multiplication by $1/4$, to
\[
 \frac{B_{p-3}(1/2)-B_{p-3}}{p-3}
 =\frac{(2^{4-p}-2)B_{p-3}}{p-3}
 \equiv-2B_{p-3}\pmod p .
\]
The Bernoulli denominators are units by von Staudt--Clausen.
It follows that $Y_3(-1)\equiv B_{p-3}/4$ and
$9A\equiv52kB_{p-3}\pmod p$.
Finally $k\equiv-(2^{p-1}-1)/p\pmod p$ by expanding the logarithm
of $2^{p-1}$. Away from $13$ the two possible zero factors give
precisely the base-two Wieferich and the Wolstenholme congruences
(the latter equivalently $\binom{2p-1}{p-1}\equiv1\pmod{p^4}$ by
\cite[Theorem~1]{AebiCairns}). At $13$ the factor $52$ vanishes.
\end{proof}

\subsection{The residue and lifting certificates}

For a residue $a\ne0,1$ let $\zeta_a$ be its Teichm\"uller representative,
$h=a/(1-a)$ and $c_m=\ell_m(a)/(a-1)$. On its disc $L$ takes all
values in $\mathbb Z_p$. Besser's reduction at the centre and the
Coleman differential equations give, for $1\le m\le4$,
\begin{equation}\label{eq:exc-local}
 Y_m(L)\equiv\sum_{j=0}^{m-1}c_{m-j}L^j/j!+hL^m/m!\pmod p.
\end{equation}
Indeed $Y_m'=Y_{m-1}$ and, for $R(x)=x/(1-x)$ and
$b_j=((x\,d/dx)^jR)(\zeta_a)$, the exact tail after the centre
terms is $\sum_{j\ge0}b_jp^jL^{j+m}/(j+m)!$. Each $b_j$ is integral.
For $p\ge7$ and $j\ge1$, $v_p((j+m)!)< (j+m)/(p-1)\le(j+4)/6$,
so this tail reduces to its $j=0$ term and converges on the whole
disc. See \cite[Propositions~2.1 and~2.3]{Besser}.
On $(1,1)$ the logarithm vanishes, so $z=\zeta_a$ and both
$c_2=c_4=0$ are necessary \cite[Theorem~3.2]{Ludtke}.

\begin{proof}[Proof of the depth-four assertion of Theorem~\ref{thm:cert}]
L\"udtke's Theorem~1.2 covers $11,1093,3511$. For each of the three
remaining primes we show that $(1,0)$ consists of $-3,-1,3,9$ and
$(1,1)$ consists of $-1$. The integral-point inclusion gives all
five component memberships. Lemma~3.1 of \cite{Ludtke} then gives
the 21 integral points in the full refined locus.

At $p=1006003$ we have $s\equiv0$ and $k,d,C$ are units. For
clarity, define
\[
 H(y)=Ad^3y_1^2y_4-Bkd^2y_1y_2y_3-Ck^3y_2^3,
 \qquad Q(X)=\ell_1(X)^2+2X\ell_2(X).
\]
Multiply \eqref{eq:exc-f4} by $d^3Y_1^2$ and substitute
\eqref{eq:exc-f2} to obtain $s^3 H(Y)=0$. Subtract
$CL^2$ times \eqref{eq:exc-f2} from $d$ times \eqref{eq:exc-f4}
and cancel $s$ to obtain
$s^2AdY_4+sBdLY_3-CkL^2Y_2=0$.
If $L\equiv0$ then $H(\boldsymbol\ell(a))=0$ by homogeneity.
Otherwise the two equations imply $Y_1\equiv Y_2\equiv0$, and
eliminating $L$ in \eqref{eq:exc-local} gives $Q(a)=0$.
At a residue with $c_1=0$, either necessary condition forces
$c_2=0$. All eight such residues have $c_2\ne0$ in the exhaustive
field calculation. The first branch leaves
\[
 3,\ 9,\ 335334,\ 670669,\ 894225,\ 1006000,\ 1006002,
\]
and the second leaves $2$. On all eight discs the depth-two reduction
$dL(c_1+hL)$ has two simple roots. Each of the sixteen classes therefore
has exactly one depth-two point. Twelve lifts have a nonzero digit of
$f_4/p^{13}$. The four other classes contain the four stated integral
points and uniqueness identifies each lift. For example, the point $2$
has $f_4(2)/p^{15}\equiv717943\pmod p$. The common centre test on
$(1,1)$ leaves only $a=-1$.

At $p=16843,2124679$ the cofactor satisfies $A\equiv0$. On every
disc the depth-two reduction $uL^2+vL+w$ has coefficients
\[
 (u,v,w)=\bigl(h(ks/2+d),\ c_1(ks+d),\ ks c_2\bigr),\qquad u\ne0.
\]
For each square discriminant, test both roots (once for zero) in
$L(s^2B Y_3(L)+CL^2Y_1(L))=0$, using \eqref{eq:exc-local}.
The enumeration gives the following counts and the first nonzero
digit of the raw determinant at $2$.
\[
\begin{array}{r|r|r|r|r|r}
p&\text{discs}&\text{depth-two classes}&\text{common}
 &\text{singular common}&f_4(2)/p^{13}\bmod p\\ \hline
16843&16841&16692&5&0&14351\\
2124679&2124677&2124776&5&0&1644662
\end{array}
\]
Evaluating $L$ at $-3,-1,2,3,9$ gives the five distinct common
classes. Hensel's lemma gives exactly one depth-two root per class.
At $2$ the root is $2$ itself, since
$\operatorname{Li}_1(2)=\operatorname{Li}_2(2)=0$; the last column
excludes it at depth four. The remaining four lifts are the stated
integral points. For both primes the simultaneous centre test
$c_2=c_4=0$ leaves just $a=-1$.
\end{proof}

The base-two degeneracy differs slightly. At $1093$ and $3511$,
$s,d$ are units, $k\equiv0$, and $A=kA'$, $B=kB'$ by the third
anchor entries. Eliminating $LY_1$ from \eqref{eq:exc-f4} gives
$s^2dA'Y_4+sdB'LY_3-CL^2Y_2=0$. Testing at the two reduced
depth-two coordinates $L=0,\ell_1(a)/a$ leaves respectively six
residues $3,9,822,972,1090,1092$ and four residues $3,9,3508,3510$.
The branches coincide at $-1$ in both cases. The field tests give
candidate counts here; the exact loci follow from \cite[Theorem~1.2]{Ludtke}.
At $11$ the same $H,Q$ test as at $1006003$ gives seven and one
residues respectively and sixteen simple depth-two classes.

\paragraph{Reproduction and precision.}
The companion source supplies the three exhaustive screens and their
reproduction commands. SageMath evaluates the four tables by Bluestein
convolution; a separate FLINT~3 checker uses a polynomial subproduct tree.
The two paths compare every table entry. At \(p=1006003\) the precision-12
calculation produces the sixteen rows of the accompanying lift table.
For $1006003$, \texttt{check\_lifts.py 1006003 26} and the same command
with $32$ independently reproduce the sixteen JSON lines using NumPy,
gmpy2, modular integration and Newton iteration. To reproduce the small
Wieferich field counts run \texttt{sage -python family\_budget.py}
with each prime as its argument. At a centre $\zeta$ the weight-$m$
depleted value is the unit moment of $x^{-m}$ for the measure with
mass $\zeta^a/(1-\zeta^{p^e})$ on $a+p^e\mathbb Z_p$.
Expansion at $a=i+pb$ through degree $r-1$ computes it modulo $p^r$,
since all higher terms lie in $p^r\mathbb Z_p$. Multiplication by
$p^m/(p^m-1)$ restores the Coleman value. In $t=(z-\zeta)/p$,
successive integrations of $d\operatorname{Li}_m=
\operatorname{Li}_{m-1}\,dz/z$ give degree-$j$ coefficients of
valuation at least $j-m\lfloor\log_p j\rfloor$. The checkers cap precision
at the coefficient and omitted-tail bounds before dividing by powers
of $p$. At the Wolstenholme primes eight depleted digits give absolute
precision $20$ for the displayed raw determinants. Lean checks the
algebraic cancellations, local reduction implications and discriminant
test. The numerical tables, Coleman input and Hensel identifications
enter through the calculations above.

\subsection{Orders three, four and six on the torsion component}

Fix $m\in\{2,4\}$ and an order $e=3$ or $4$. For $p\equiv1\pmod e$
choose a primitive $e$th root $\eta\in\mathbb Z_p$. Let $\chi$ be
the primitive odd character of conductor $e$ and $\omega$ the
Teichm\"uller character. In the interpolation convention
$L_p(1-j,\psi)=(1-(\psi\omega^{-j})(p)p^{j-1})L(1-j,\psi\omega^{-j})$
with the primitive character in the Euler factor, we have
\[
 2M_m(\eta):=2(1-p^{-m})\operatorname{Li}_m(\eta)
 =(\eta-\eta^{-1})L_p(m,\chi\omega^{1-m}).
\]
Indeed $\eta^a-\eta^{-a}=(\eta-\eta^{-1})\chi(a)$. For
$n_r=p^{r-1}(p-1)-m$, the difference of the full power moments at
$n_r$ equals $(\eta-\eta^{-1})L(-n_r,\chi)$ by the Bernoulli
generating series. Removing multiples of $p$ contributes
$1-p^{n_r}$, and interpolation as $-n_r\to m$ gives the formula.
The inverse-root moments have opposite signs by even-weight inversion
\cite[\S\S VI--VII]{Coleman}. For a primitive sixth root $\xi$,
$-\xi=(\xi^2)^{-1}$. Duplication and inversion give
$\operatorname{Li}_m(\xi)=(1+2^{1-m})\operatorname{Li}_m(\xi^2)$.
The factors $3/2,9/8$ are units for the primes in this section.

For $e=3$, put $h=p-m$. Applying the Bernoulli formula for
$L(1-h,\chi)$ and reflecting $B_h(2/3)=-B_h(1/3)$ gives
\[
 M_m(\eta)\equiv
 -\frac{(\eta-\eta^{-1})3^{h-1}}{h}B_h(1/3)\pmod p .
\]
Sun reports one prime $3<p<2\cdot10^7$ with
$B_{p-2}(1/3)\equiv0$, namely $205129$ \cite[Remark~17]{Sun}.
At its cubic root of residue $81187$ one finds
$M_2/p\equiv170025$ and $M_4\equiv30238\pmod p$.
Thus weight two excludes orders three and six on $(1,1)$ throughout
$7\le p<2\cdot10^7$, using a second digit at that prime. At nonsplit
primes the roots do not belong to $\mathbb Q_p$.

For $e=4$ let $i^2=-1$, $n=p-1-m$ and define signed Euler numbers by
$\operatorname{sech}(t)=\sum E_jt^j/j!$. Differentiating
$(1-X^p)/(1-X)=\sum_{j=0}^{p-1}X^j$ with $(X\,d/dX)^n$ in
characteristic $p$ gives
\[
 \ell_m(i)=(1-i)\frac{i}{2}E_n .
\]
Indeed the $n$th logarithmic derivative of $(1-X)^{-1}$ at $i$
is the $n$th derivative at zero of
$(1-i e^t)^{-1}=(1-\tanh t+i\operatorname{sech} t)/2$,
and $n$ is positive even. Me\v{s}trovi\'c's exhaustive computation
finds just three primes below $10^7$ dividing $E_{p-3}$
\cite{Mestrovic}. For the indicated roots the following table records
the next weight and the next digit of the original weight.
\[
\begin{array}{r|r|r|r|r}
p&i\bmod p&\ell_2(i)&\ell_4(i)&M_2(i)/p\bmod p\\ \hline
149&44&0&143&85\\
241&64&0&105&103\\
2946901&722077&0&1422053&1795894
\end{array}
\]
Here the $\ell_m$ columns are finite-field values. For either
conductor the moment column is reproducible modulo $p^2$ by summing
$\eta^a a^{-m}/(1-\eta)$ over $1\le a<p^2$, $p\nmid a$.
In $a=u+pb$, the two inner sums are
$\sum_{b=0}^{p-1}\eta^b=1$ and
$\sum_{b=0}^{p-1}b\eta^b=(p-1)\eta/(\eta-1)$; expanding the inverse
power to first order leaves $p-1$ outer summands. Inversion handles
the conjugate root. This excludes primitive fourth roots on $(1,1)$
for $7\le p<10^7$. Both range conclusions use the cited searches
and direct calculations at the listed exceptions. They concern these
specified torsion orders; the distribution of other orders and the
cross-prime density of simultaneous zeros remain open here.

\section{Limits of the estimates and a projective comparison}
\subsection{The dimension ceiling in four weights}
Let \(J_{p,t}\) evaluate \(P(X,y_1,\ldots,y_4)\) at the four finite
polylogarithms modulo \(X^p\), with \(\deg_X P\le t^2\) and
\(\deg_y P\le t\) in total. Its source has dimension
\(D=(t^2+1)\binom{t+4}{4}\sim t^6/24\). An injectivity theorem up
to \(t=\lfloor cp^{1/6}\rfloor\) for a fixed small \(c>0\) would give
\(O(p^{5/6})\) in Lemma~\ref{lem:mixed-disc-count} for \(n=4\), by
leaving the auxiliary polynomial and multiplicity steps unchanged.
The example in Theorem~\ref{thm:zero} explains why an improvement of this
kind must exploit the particular system of finite polylogarithms.

For \(t=1,2,3\) one can prove maximal rank for every prime by an integer
calculation. On coefficients of index at most \(D\), the Taylor coefficients
of the characteristic-zero series \(\sum_{j\ge1}X^j/j^i\) specialise
to the finite series for every \(p>D\). The pure \(X\)-power columns give
an identity block on rows \(0,\ldots,t^2\). Complete the block in two ways
by using either the next \(D-t^2-1\) rows or the same rows shifted up
by one index. Clear denominators in each remaining row. The clearing
factors are units for \(p>D\). If both integer minors vanished modulo such a
prime it would divide their gcd. The following table gives the bit lengths
of the minors, their gcd after removing all factors at most \(D\), and
the number of smaller primes checked directly for surjectivity.
\begin{center}
\begin{tabular}{rrrrr}
\toprule
\(t\)&\(D\)&bit lengths&remaining gcd&primes \(p\le D\)\\
\midrule
1&10&37,\ 37&1&4\\
2&75&9636,\ 9869&1&21\\
3&350&332544,\ 334400&1&70\\
\bottomrule
\end{tabular}
\end{center}
Thus \(\operatorname{rank}J_{p,t}=\min(D,p)\) for these three values of \(t\).
The second minor matters already at \(t=2,p=109\), where the first 75
rows have rank 74. In boundary tests up to \(t=7\), the calculation gives
ranks \(16493\) and \(16500\) at \(p=16493\) and \(16519\) respectively
when \(D=16500\). A bound for growing \(t\) remains open.

\subsection{Density one at depth four}
For \(S=\{2,3\}\), equality with the 21 integral points for a density-one
set of auxiliary primes would require excluding extra common zeros on
both refinement components. Even on \((1,1)\), where the logarithm
restricts the points to roots of unity, two difficulties remain. The
possible orders grow with \(p-1\). At a fixed order three or four the
conditions at weights two and four evaluate different branches
\(\chi\omega^{-1}\) and \(\chi\omega^{-3}\) of the \(p\)-adic \(L\)-function.
The calculation at either conductor in Section~\ref{sec:exceptional}
rules out individual primes by finding a nonzero digit. It gives no
density estimate for simultaneous vanishing as \(p\) varies.
Betts, Kumpitsch and L\"udtke isolate \(-1\) using a weight depending on
\(p\) \cite[Proposition~5.15]{BKL}. The question here fixes the depth.

\subsection{A genus-three calculation and the Cartier operator}
Consider the smooth projective model \(C\) of
\[
 y^2=f(x)=4x^7+9x^6-8x^5-36x^4-16x^3+32x^2+32x+8.
\]
Its Jacobian \(J\) has rank one. The divisor \(D=(-1,-1)-\infty\) has
infinite order \cite[Example~4.1]{GenusThree}. At a good odd prime we
count the discs meeting the inverse image of \(\mathbb Q_p\log D\)
under \(P\mapsto\log[P-\infty]\). This definition uses the rational
Mordell--Weil span and requires no identification with a Selmer group.

Away from branch points choose a lift \(Q\) with local coordinate
\(x=x(Q)+pw\). Project the integrals of \(x^i\,dx/(2y)\), \(0\le i\le2\),
onto an integral basis of the two-dimensional annihilator of \(\log D\).
Write \(c_1,c_2\) for the two projected constants and \(u_1,u_2\) for their
tangent coefficients modulo \(p\). A lift requires \(c_i\in p\mathbb Z_p\)
and
\[
                 c_i/p+u_iw\equiv 0 \pmod p,\qquad i=1,2.
\]
For a nonzero slope vector, compatibility amounts to
\((c_1/p)u_2-(c_2/p)u_1\equiv0\pmod p\). When both slopes vanish, both
reduced constants must vanish. A change of lift translates the constant
vector along the slope vector. Stoll also uses the direction of a
normalised abelian logarithm, combining it at \(p=2\) with a global
Selmer-group test for rational points \cite[Proposition~2.8]{Stoll}.

Table~\ref{tab:genus} gives the calculation for all 45 odd primes below
200. For \(n=\#J(\mathbb F_p)\), compute \([n](Q-\infty)\) in the
formal neighbourhood of the identity. Translating by a regular
degree-three divisor allows Cantor arithmetic on a chart with simple
support. Integrating between the two sets of nearby support points and
taking the trace gives \(p^{-1}\log([n](Q-\infty))\bmod p^3\).
For \(p\ge5\) the omitted terms of degree \(j\ge4\) have valuation at least
\(j-v_p(j)\ge4\) before dividing by \(p\). Testing all admissible local
digits gives the upper count. Lower witnesses are the Weierstrass
points and the available pairs above \(x=-1,1,0\). The former two
pairs are rational. For \(x=0\) a torsion witness follows from
\[
 \begin{aligned}
 B &=x^2+4x+2,\\
 A &=-x^6-16x^5-36x^4+72x^2+64x+16,&
 A^2-8fB^2&=x^{12}.
 \end{aligned}
\]
At \(P=(0,2\sqrt2)\) the function \(A-2\sqrt2\,yB\) has a zero
of order 12 and a unique pole at infinity. Its divisor proves that
\(12[P-\infty]=0\). For \(p=3,7\) there is a witness on every disc.
At the remaining primes the upper and lower counts coincide.
In particular, at \(p=41\) successive precision levels reduce the
upper count from \(40\) to \(12\) and then \(10\).
\begin{table}[htbp]
\centering
\caption{Classical abelian Chabauty on \(C\). In each group \(n\) counts
all residue discs and \(o\) counts the occupied discs.}\label{tab:genus}
\begin{tabular}{rrr@{\qquad}rrr@{\qquad}rrr}
\toprule
\(p\)&\(n\)&\(o\)&\(p\)&\(n\)&\(o\)&\(p\)&\(n\)&\(o\)\\
\midrule
3&5&5&59&61&5&131&161&5\\
5&6&4&61&62&8&137&140&10\\
7&10&10&67&69&5&139&138&8\\
11&16&6&71&79&9&149&164&6\\
13&15&5&73&72&10&151&140&12\\
17&15&11&79&75&11&157&152&6\\
19&28&6&83&72&6&163&165&7\\
23&30&10&89&81&9&167&139&9\\
29&30&6&97&103&9&173&176&6\\
31&29&9&101&106&6&179&183&5\\
37&39&5&103&103&9&181&185&5\\
41&40&10&107&96&6&191&201&11\\
43&47&7&109&109&5&193&208&10\\
47&45&9&113&124&10&197&211&5\\
53&62&6&127&145&9&199&184&10\\
\bottomrule
\end{tabular}
\end{table}

In the holomorphic basis the Cartier--Manin matrix has entries
\([x^{pi-j}]f^{(p-1)/2}\), \(1\le i,j\le3\), up to the transpose
convention \cite{AchterHowe}. Its rank is three at 43 of these primes
and two at 5 and 13. Since Cartier annihilates exact differentials, a
nonzero reduced annihilating differential at an ordinary prime cannot
have a rational primitive on the special fibre. The tangent slopes in
the test are algebraic; recovering the constants across discs calls for
further information about the lifts. Obtaining a bounded-degree relation
among those constants as \(p\) varies remains open.

\paragraph{Formalisation and reproducibility.}
The companion source is available on GitHub.
\begin{quote}
\url{https://github.com/prajaybasu/chabauty-kim-disc-bounds}
\end{quote}
It contains a single Lean package. Its logarithmic modules check the last-pivot row, adjugate denominator, bounded-height
horizontal classification, jet differentiation, and the dimension and
kernel calculations used for the lower bound. The choice of the cyclic
rows, their degree growth and the realisation by iterated integrals are
given in Section~\ref{sec:zero}. The Selmer and disc modules verify
bounded-degree elimination, primitive reduction and inverse translation.
Localisation, the Coleman comparison and the auxiliary-polynomial count
are supplied by the cited results and Section~\ref{sec:discs}.
The exceptional-prime modules verify the local tail and finite-field
implications with the stated analytic inputs. The exact tables are
checked by arithmetic programs.

\begin{samepage}
To reproduce the rank computation, run the following command.\par\nopagebreak
\begin{verbatim}
sage -python mixed_rank.py --certify 1 2 3
\end{verbatim}
\end{samepage}
To reproduce the 45-prime census and its witnesses, run this command in
the same directory.
\begin{verbatim}
sage -python genus_three_jacobian.py
\end{verbatim}

\bibliographystyle{plain}
\bibliography{references}

\noindent Email: \href{mailto:prajay@prajaybasu.com}{\texttt{prajay@prajaybasu.com}}
\end{document}
